\documentclass[11pt]{article}
\usepackage{amsmath, amssymb}
\usepackage{geometry}
\usepackage{array}
\usepackage{booktabs}
\geometry{margin=0.65in}
\setlength{\parindent}{0pt}
\setlength{\parskip}{0.3em}

\begin{document}
\begin{center}
{\LARGE Algebra II Topic Cheatsheet}\\[0.4em]
{\large Comprehensive Review}
\end{center}

\section*{Polynomials and Factoring}
\[
a^2-b^2=(a-b)(a+b)
\]
\[
a^2\pm 2ab+b^2=(a\pm b)^2
\]
\[
a^3-b^3=(a-b)(a^2+ab+b^2),\qquad
a^3+b^3=(a+b)(a^2-ab+b^2)
\]
\begin{itemize}
  \item Factor out the GCF first.
  \item Use grouping for four terms.
  \item For $ax^4+bx^2+c$, substitute $u=x^2$.
  \item For quartics, look for symmetric factors of the form $(x^2+px+q)(x^2-px+r)$.
\end{itemize}

\section*{Rational Root Theorem and Division}
If
\[
f(x)=a_nx^n+\cdots+a_1x+a_0
\]
has a rational root $\dfrac{p}{q}$ in lowest terms, then
\[
p\mid a_0,\qquad q\mid a_n
\]
Possible rational roots:
\[
\pm \frac{\text{factors of constant term}}{\text{factors of leading coefficient}}
\]
Remainder Theorem:
\[
f(r)=\text{remainder on division by }x-r
\]
Factor Theorem:
\[
f(r)=0 \iff x-r\text{ is a factor}
\]
Synthetic division by $x-r$ uses the number $r$.

\section*{Zeros and Graph Behavior}
\begin{itemize}
  \item Real zeros give $x$-intercepts.
  \item Even multiplicity: touches and turns.
  \item Odd multiplicity: crosses.
  \item Complex roots of polynomials with real coefficients come in conjugate pairs.
  \item End behavior comes from degree and leading coefficient.
\end{itemize}
\[
\begin{array}{ll}
\text{even degree, positive LC} & \text{both ends up} \\
\text{even degree, negative LC} & \text{both ends down} \\
\text{odd degree, positive LC} & \text{left down, right up} \\
\text{odd degree, negative LC} & \text{left up, right down}
\end{array}
\]

\section*{Rational Functions}
For
\[
f(x)=\frac{p(x)}{q(x)}
\]
\begin{itemize}
  \item Domain excludes zeros of $q(x)$.
  \item Vertical asymptotes occur at non-canceled zeros of $q(x)$.
  \item Holes occur when a factor cancels.
  \item Horizontal asymptotes:
  \[
  \deg p<\deg q \Rightarrow y=0
  \]
  \[
  \deg p=\deg q \Rightarrow y=\frac{\text{leading coeff of }p}{\text{leading coeff of }q}
  \]
  \[
  \deg p>\deg q \Rightarrow \text{use long division}
  \]
  \item Slant asymptote occurs when $\deg p=\deg q+1$.
\end{itemize}
Intercepts:
\[
x\text{-intercepts from numerator zeros not canceled},\qquad y\text{-intercept }f(0)
\]

\section*{Logarithms and Exponentials}
\[
y=b^x \iff x=\log_b y
\]
\[
\log_b b^x=x,\qquad b^{\log_b x}=x
\]
with
\[
b>0,\quad b\ne 1,\quad x>0
\]
\[
\log_b(MN)=\log_bM+\log_bN
\]
\[
\log_b\!\left(\frac{M}{N}\right)=\log_bM-\log_bN
\]
\[
\log_b(M^p)=p\log_bM
\]
\[
\log_b a=\frac{\log a}{\log b}=\frac{\ln a}{\ln b}
\]
Solving templates:
\[
b^{f(x)}=b^{g(x)} \Rightarrow f(x)=g(x)
\]
\[
\log_b(f(x))=\log_b(g(x)) \Rightarrow f(x)=g(x)
\]
Check domain: log arguments must be positive.

Growth and decay:
\[
y=ab^t,\qquad y=ae^{kt}
\]
\[
k>0 \Rightarrow \text{growth},\qquad k<0 \Rightarrow \text{decay}
\]
\[
\text{doubling time}=\frac{\ln 2}{k},\qquad
\text{half-life}=\frac{\ln(1/2)}{k}
\]

\section*{Quadratics and Conics}
Vertex form:
\[
y=a(x-h)^2+k
\]
Axis of symmetry:
\[
x=h
\]
Vertex:
\[
(h,k)
\]
Maximum or minimum depends on the sign of $a$.

Conic identification:
\[
Ax^2+Bxy+Cy^2+Dx+Ey+F=0
\]
\begin{itemize}
  \item circle: $A=C$, $B=0$
  \item parabola: one squared term
  \item ellipse: same sign, unequal coefficients
  \item hyperbola: squared terms with opposite signs
\end{itemize}

Standard forms:
\[
(x-h)^2+(y-k)^2=r^2
\]
\[
(x-h)^2=4p(y-k),\qquad (y-k)^2=4p(x-h)
\]
\[
\frac{(x-h)^2}{a^2}+\frac{(y-k)^2}{b^2}=1
\]
\[
\frac{(x-h)^2}{a^2}-\frac{(y-k)^2}{b^2}=1,\qquad
\frac{(y-k)^2}{a^2}-\frac{(x-h)^2}{b^2}=1
\]
For hyperbolas:
\[
c^2=a^2+b^2,\qquad e=\frac ca,\qquad \text{asymptotes } y-k=\pm \frac{b}{a}(x-h)
\]
or
\[
y-k=\pm \frac{a}{b}(x-h)
\]

\section*{Trigonometry}
\[
180^\circ=\pi,\qquad 360^\circ=2\pi
\]
\[
\sin\theta=\frac{\text{opp}}{\text{hyp}},\qquad
\cos\theta=\frac{\text{adj}}{\text{hyp}},\qquad
\tan\theta=\frac{\text{opp}}{\text{adj}}
\]
\[
\csc\theta=\frac1{\sin\theta},\qquad
\sec\theta=\frac1{\cos\theta},\qquad
\cot\theta=\frac1{\tan\theta}
\]
Core identities:
\[
\sin^2\theta+\cos^2\theta=1
\]
\[
1+\tan^2\theta=\sec^2\theta,\qquad
1+\cot^2\theta=\csc^2\theta
\]
Special values:
\[
\sin\left(\frac{\pi}{6},\frac{\pi}{4},\frac{\pi}{3}\right)=\left(\frac12,\frac{\sqrt2}{2},\frac{\sqrt3}{2}\right)
\]
\[
\cos\left(\frac{\pi}{6},\frac{\pi}{4},\frac{\pi}{3}\right)=\left(\frac{\sqrt3}{2},\frac{\sqrt2}{2},\frac12\right)
\]
\[
\tan\left(\frac{\pi}{6},\frac{\pi}{4},\frac{\pi}{3}\right)=\left(\frac{\sqrt3}{3},1,\sqrt3\right)
\]
Law of Sines:
\[
\frac{a}{\sin A}=\frac{b}{\sin B}=\frac{c}{\sin C}
\]
Law of Cosines:
\[
c^2=a^2+b^2-2ab\cos C
\]
Triangle area:
\[
K=\frac12 ab\sin C
\]

\section*{Sequences and Series}
Arithmetic sequence:
\[
a_n=a_1+(n-1)d
\]
\[
S_n=\frac{n(a_1+a_n)}{2}
\]
Geometric sequence:
\[
a_n=a_1r^{n-1}
\]
\[
S_n=\frac{a_1(1-r^n)}{1-r}\qquad (r\ne 1)
\]
Infinite geometric series:
\[
S_\infty=\frac{a_1}{1-r}\qquad (|r|<1)
\]
Sigma notation:
\[
\sum_{k=1}^{n} a_k
\]
\[
\text{factorials: }n!=n(n-1)\cdots 1
\]
\[
{}_nP_r=\frac{n!}{(n-r)!},\qquad
{}_nC_r=\frac{n!}{r!(n-r)!}
\]

\section*{Complex Numbers and Polar}
\[
i^2=-1,\qquad \sqrt{-a}=i\sqrt a
\]
If
\[
z=a+bi,
\]
then its conjugate is
\[
a-bi
\]
and
\[
(a+bi)(a-bi)=a^2+b^2
\]
Polar form:
\[
z=r(\cos\theta+i\sin\theta)=r\,\text{cis}\,\theta
\]
Polar to Cartesian:
\[
x=r\cos\theta,\qquad y=r\sin\theta
\]
Cartesian to polar:
\[
r^2=x^2+y^2,\qquad \tan\theta=\frac{y}{x}
\]
De Moivre:
\[
\big[r(\cos\theta+i\sin\theta)\big]^n=r^n(\cos n\theta+i\sin n\theta)
\]
Nth roots of $r(\cos\theta+i\sin\theta)$:
\[
z_k=r^{1/n}\left(\cos\frac{\theta+2\pi k}{n}+i\sin\frac{\theta+2\pi k}{n}\right),\quad k=0,1,\dots,n-1
\]

\section*{Quick Algebra II Checks}
\begin{itemize}
  \item When solving, check for extraneous roots from logs and rational equations.
  \item Use exact values whenever possible.
  \item State rounding only when asked.
  \item For rational equations, multiply by the LCD only after listing domain restrictions.
\end{itemize}

\end{document}
